% ==============================================================================
% SEÇÃO 6: CRITÉRIOS DE ACEITE E HOMOLOGAÇÃO FINAL
% Autor: Danilo Belém Oliveira (Issue #10)
% 
% REGRA DE OURO ANTI-REDUNDÂNCIA:
% Esta seção funciona como o checklist formal de homologação e auditoria 
% para a professora e banca examinadora. Formato predominantemente tabular.
% ==============================================================================

\section{Critérios de Aceite e Homologação Final}
\label{sec:criterios_aceite}

Este capítulo formaliza o processo de homologação do sistema frente aos requisitos pedagógicos da disciplina, apresentando os critérios quantitativos e verificações práticas que comprovam o sucesso da engenharia e da modelagem desenvolvidas.

% ------------------------------------------------------------------------------
% 6.1 HOMOLOGAÇÃO DAS 5 HIPÓTESES DA EDA
% ------------------------------------------------------------------------------
\subsection{Homologação das Hipóteses Mercadológicas da EDA}
\label{sec:homologacao_hipoteses}

% [PLACEHOLDER - DANILO: Tabela de verificação empírica de H1 a H5]
% - Hipótese x Afirmação Teórica x Resultado Empírico Auditado x Status (Confirmada/Refutada);
% - Uso rigoroso de medianas e linguagem epistêmica prudente.
\noindent\textit{[Espaço reservado para a tabela de homologação das 5 hipóteses da EDA por Danilo Belém -- Issue \#10.]}

% ------------------------------------------------------------------------------
% 6.2 HOMOLOGAÇÃO PRÁTICA DAS 5 USER STORIES
% ------------------------------------------------------------------------------
\subsection{Homologação Prática das User Stories no Dashboard}
\label{sec:homologacao_us}

% [PLACEHOLDER - DANILO: Tabela de verificação de US1 a US5]
% - História x Tela no Streamlit x Condição de Teste Prático x Status de Aceite (100% Atendido);
% - Referenciar a Matriz Geral de Rastreabilidade da Seção 4 (Tabela~\ref{tab:matriz_rastreabilidade}).
\noindent\textit{[Espaço reservado para a tabela de verificação das 5 User Stories no dashboard por Danilo Belém -- Issue \#10.]}

% ------------------------------------------------------------------------------
% 6.3 OS 9 CRITÉRIOS FORMAIS DE AUDITORIA ALGORÍTMICA (C1 A C9)
% ------------------------------------------------------------------------------
\subsection{Critérios Formais de Auditoria e Validação Algorítmica}
\label{sec:criterios_auditoria}

% [PLACEHOLDER - DANILO: Tabela formal dos critérios C1 a C9 de AUDITORIA_METRICAS.md]
% - C1: Redução de RMSE em relação à Média Global;
% - C2: Superioridade do SVD frente ao KNN (Wilcoxon p < 0,05);
% - C3: Precision@10 e NDCG@10 do SVD superando baselines;
% - C4: Explicabilidade aditiva vetorial sem resíduo;
% - C5: Imunidade a cold start na CBF (similaridade >= 0,5);
% - C6: Governança e ausência de PII sob a LGPD;
% - C7: Desempate de catálogo pelo escore bruto contínuo;
% - C8: Piso teórico do oráculo (RMSE = 0,5214 com SVD em 1,20x);
% - C9: Avaliação formal da CBF sob o protocolo de ranking (P@10 = 0,882).
\noindent\textit{[Espaço reservado para a tabela dos 9 critérios formais de auditoria algorítmica por Danilo Belém -- Issue \#10.]}
